%%%%%%%%%%%%%%%%%%%%%%%%%%%%%%%%%%%%%%%%%
% Important note:
% This template must be compiled with XeLaTeX, the below lines will ensure this
%!TEX TS-program = xelatex
%!TEX encoding = UTF-8 Unicode
%
%%%%%%%%%%%%%%%%%%%%%%%%%%%%%%%%%%%%%%%%%

%----------------------------------------------------------------------------------------
%	PACKAGES AND OTHER DOCUMENT CONFIGURATIONS
%----------------------------------------------------------------------------------------

\documentclass[11pt, letterpaper]{awesome-cv} % A4 paper size by default, use 'letterpaper' for US letter

\geometry{left=2cm, top=1.5cm, right=2cm, bottom=2cm, footskip=.5cm} % Configure page margins with geometry
 
\fontdir[fonts/] % Specify the location of the included fonts

% Color for highlights
%\colorlet{awesome}{awesome-skyblue} % Default colors include: awesome-emerald, awesome-skyblue, awesome-red, awesome-pink, awesome-orange, awesome-nephritis, awesome-concrete, awesome-darknight
%\definecolor{awesome}{HTML}{CA63A8} % Uncomment if you would like to specify your own color
\definecolor{awesome}{HTML}{0BA89A} % UCSB sea green

% Colors for text - uncomment and modify
%\definecolor{darktext}{HTML}{000000}
\definecolor{text}{HTML}{000000}
\definecolor{graytext}{HTML}{000000}
%\definecolor{lighttext}{HTML}{000000}

\renewcommand{\acvHeaderSocialSep}{\quad\textbar\quad} % If you would like to change the social information separator from a pipe (|) to something else

%----------------------------------------------------------------------------------------
%	PERSONAL INFORMATION
%	Comment any of the lines below if they are not required
%----------------------------------------------------------------------------------------

\name{Jason C.}{Wei}
\address{Pleasanton, California, USA}
\mobile{(925) 352-7600}
\github{Jayweiii}
\homepage{jweii.com}
\email{jwei@ucsb.edu}
\linkedin{jwei}
%\skype{skypeid}
%\stackoverflow{SOid}{SOname}
%\twitter{@twit}
%\linkedin{linkedin name}
%\reddit{reddit account}
%\xing{xing name}
%\extrainfo{test} % Other text you want to include on this line

\position{Mechanical Engineer{\enskip\cdotp\enskip}UCSB} % Your expertise/fields
\quote{Rising third-year MechE at UCSB with extensive hands-on experience in \\ 3D printing, Formula SAE EV powertrain design, phase-change thermofluidics, and CNC Milling.} % A quote or statement

%----------------------------------------------------------------------------------------
%	RECIPIENT/POSITION/LETTER INFORMATION
%	All of the below lines must be filled out
%----------------------------------------------------------------------------------------

\recipient{Hiring Manager}{Seek Thermal\\Santa Barbara, CA} % The company being applied to

\letterdate{\today} % The date on the letter, default is the date of compilation

\lettertitle{Mechanical Engineering Intern} % The title of the letter

\letteropening{Dear Seek Thermal,} % How the letter is opened

\letterclosing{Kind regards,} % How the letter is closed

% \letterenclosure[Attached]{Resume} % Any enclosures with the letter

\makecvfooter{\today}{Jason C. Wei~~~·~~~Cover Letter}{} % Specify the letter footer with 3 arguments: (<left>, <center>, <right>), leave any of these blank if they are not needed
  
%----------------------------------------------------------------------------------------

\begin{document}

\makecvheader % Print the header

\makelettertitle % Print the title

%----------------------------------------------------------------------------------------
%	LETTER CONTENT
%----------------------------------------------------------------------------------------

\begin{cvletter}

%------------------------------------------------

I am delighted to submit my application for the Mechanical Engineering Intern position at Seek Thermal. As a rising third-year Mechanical Engineering student at the University of California - Santa Barbara with extensive hands-on experience in, SolidWorks, CNC Manufacturing, and Formula SAE EV powertrain design,  I am confident that my skills and experiences make me an excellent candidate for this internship.

In my role as an undergraduate researcher at Prof. Yangying's Thermofluid Energy Research Laboratory, I have worked very closely with a PhD student and manufactured a flow boiling testing apparatus using 3D printing and CNC milling, which went up to the International Space Station in November 2022 as part of a NASA mission. Additionally, as a personal research project under the mentorship of Yangying, I have been exploring combining 3D printed internal geometries, such as gyroids, with gas-liquid double phase change, to increase the efficiency of heat exchangers. 

At Gaucho Racing, I lead 60 engineers in the design and manufacturing of our high voltage battery pack, drivetrain, and mounting solutions. To tackle our design criteria most cost and time effectively, I utilize various manufacturing methods including CNC, welding, 3D printing, laser cutting, soldering, and resin bonding. 

As a Workshop Wizard and 3D Printer Tool Lead at the CNSI Microfluidics lab and Innovation Workshop, I have manufactured and designed over 100 researcher experimental setups. Additionally, I manage, repair, and improve processes for our Haas Super Mini Mill and all types of our 3D printers, including FDM, SLA, DLP, and PolyJet. These experiences have honed my understanding of the different design and material considerations of additive manufacturing, subtractive manufacturing (CNC \& manual machining), chemical post-processes (such as acetone vapor smoothing), and even combining these different manufacturing methods in creating a single part.

I am excited to bring the research, design, and manufacturing skills that I have accrued through hands-on engineering experiences to Seek Thermal. Thank you for considering my application. I look forward to discussing this internship opportunity and my qualifications in further detail.

%------------------------------------------------

\end{cvletter}

%----------------------------------------------------------------------------------------

\makeletterclosing % Print the signature and enclosures

\end{document}